\documentclass[11pt,a4paper]{article}

% ── Packages ──────────────────────────────────────────────────────────────────
\usepackage[utf8]{inputenc}
\usepackage[T1]{fontenc}
\usepackage[margin=2.5cm]{geometry}
\usepackage{hyperref}
\usepackage{xcolor}
\usepackage{listings}
\usepackage{enumitem}
\usepackage{booktabs}
\usepackage{longtable}
\usepackage{amsmath,amssymb}
\usepackage{fancyhdr}
\usepackage{titlesec}

% ── Colors ────────────────────────────────────────────────────────────────────
\definecolor{isls-blue}{HTML}{1a5276}
\definecolor{isls-orange}{HTML}{e67e22}
\definecolor{isls-green}{HTML}{27ae60}
\definecolor{isls-red}{HTML}{c0392b}
\definecolor{isls-gray}{HTML}{7f8c8d}
\definecolor{code-bg}{HTML}{f7f9fb}

% ── Listings ──────────────────────────────────────────────────────────────────
\lstdefinestyle{rust}{
  language=C,
  basicstyle=\ttfamily\small,
  keywordstyle=\color{isls-blue}\bfseries,
  commentstyle=\color{isls-gray}\itshape,
  stringstyle=\color{isls-green},
  backgroundcolor=\color{code-bg},
  frame=single,
  rulecolor=\color{isls-gray!30},
  breaklines=true,
  showstringspaces=false,
  tabsize=4,
  morekeywords={fn,let,mut,pub,struct,impl,enum,use,mod,self,Self,
    crate,super,where,trait,for,in,if,else,match,return,async,await,
    Ok,Err,Some,None,Vec,String,HashMap,HashSet,PathBuf,Result,Option},
}
\lstdefinestyle{bash}{
  language=bash,
  basicstyle=\ttfamily\small,
  backgroundcolor=\color{code-bg},
  frame=single,
  rulecolor=\color{isls-gray!30},
  breaklines=true,
  showstringspaces=false,
}
\lstdefinestyle{toml}{
  basicstyle=\ttfamily\small,
  backgroundcolor=\color{code-bg},
  frame=single,
  rulecolor=\color{isls-gray!30},
  breaklines=true,
  showstringspaces=false,
  commentstyle=\color{isls-gray}\itshape,
  morecomment=[l]{\#},
}

% ── Headers ───────────────────────────────────────────────────────────────────
\pagestyle{fancy}
\fancyhf{}
\fancyhead[L]{\small\textcolor{isls-gray}{ISLS D6 Spec --- M\"obius-Umst\"ulpung}}
\fancyhead[R]{\small\textcolor{isls-gray}{v1.0.0 --- \today}}
\fancyfoot[C]{\thepage}

\titleformat{\section}
  {\Large\bfseries\color{isls-blue}}{\thesection}{1em}{}
\titleformat{\subsection}
  {\large\bfseries\color{isls-blue!80}}{\thesubsection}{1em}{}
\titleformat{\subsubsection}
  {\normalsize\bfseries\color{isls-blue!60}}{\thesubsubsection}{1em}{}

\hypersetup{
  colorlinks=true,
  linkcolor=isls-blue,
  urlcolor=isls-orange,
  citecolor=isls-green,
}

% ══════════════════════════════════════════════════════════════════════════════
\begin{document}

% ── Titlepage ─────────────────────────────────────────────────────────────────
\begin{titlepage}
\centering
\vspace*{3cm}

{\Huge\bfseries\color{isls-blue} ISLS D6 Specification}\\[0.5cm]
{\LARGE\color{isls-orange} M\"obius-Umst\"ulpung:\\Der Generator generiert sich selbst}\\[2cm]

{\large
\begin{tabular}{rl}
\textbf{System:}    & ISLS --- Intelligent Semantic Ledger Substrate \\
\textbf{Version:}   & v4.0 (D6) \\
\textbf{Author:}    & Sebastian Klemm \\
\textbf{Date:}      & \today \\
\textbf{Status:}    & Specification \\
\textbf{Depends:}   & D5 (complete), all 10 crates \\
\textbf{Branch:}    & \texttt{claude/implement-isls-d6-spec} \\
\end{tabular}
}

\vfill

{\small\color{isls-gray}
Mithras Substructure University\\
Repository: \texttt{https://github.com/LashSesh/graphity}\\
License: MIT
}
\end{titlepage}

% ── TOC ───────────────────────────────────────────────────────────────────────
\tableofcontents
\newpage


% ══════════════════════════════════════════════════════════════════════════════
\section{Executive Summary}
% ══════════════════════════════════════════════════════════════════════════════

D1--D5 proved that ISLS can generate arbitrary CRUD applications from
natural language. D6 asks the decisive question: \textbf{can the generator
generate an application whose domain IS the generator itself?}

This is the M\"obius-Umst\"ulpung --- the topological inversion where
inside becomes outside. The generator (ISLS) produces a web application
(ISLS Studio) that manages the generator's own data: norms, candidates,
generation runs, scrape jobs. The generated application is not a toy ---
it replaces the hand-built \texttt{isls-gateway} Axum prototype with a
full Actix-Web/PostgreSQL/Docker production system.

In Dual Fabric terms: the Trumpet emits a signal, and that signal
\emph{is a mirror of the Trumpet}. The emitted structure (ISLS Studio)
reflects the emitter (ISLS) --- not as a copy of source code, but as a
self-describing domain model rendered through the generation pipeline.

\subsection{The Conceptual Leap}

Before D6, ISLS generates applications for \emph{external} domains
(warehouse, hotel, clinic, gym). After D6, ISLS generates an application
for \emph{its own} domain. The entities are not Products and Orders ---
they are Norms, NormCandidates, Generations, and ScrapeJobs.

This proves that the 48 structural generators, the HDAG pipeline, and
the LLM-driven query layer are genuinely domain-agnostic. If they can
produce a working application for the ISLS domain itself, they can
produce one for any domain.

\subsection{What D6 Achieves}

\begin{enumerate}[nosep]
\item \textbf{Self-Description TOML:} ISLS's core data types expressed
      as a standard \texttt{[[entities]]} specification.
\item \textbf{Meta-Forge Command:} \texttt{isls forge-self} generates
      ISLS Studio from the self-description.
\item \textbf{Norms Import Seed:} The generated database migration
      includes seed data for the 24 builtin norms, bootstrapping
      the Studio with a populated catalog.
\item \textbf{Compiled, Deployed, Operational:} The generated Studio
      passes \texttt{cargo build}, \texttt{docker-compose up}, login,
      and full CRUD on all entities.
\end{enumerate}

\subsection{What D6 is NOT}

\begin{itemize}[nosep]
\item \textbf{Not source-code reproduction.} The generated Studio is
      a CRUD web app that \emph{manages} ISLS data. It does not
      contain the forge pipeline, the HDAG builder, or the structural
      generators. That would be D8 territory.
\item \textbf{Not a replacement for the CLI.} The CLI (\texttt{isls
      forge-chat}, \texttt{isls scrape}, etc.) remains the primary
      operator interface. The Studio is a complementary web dashboard.
\item \textbf{Not the gateway refactor.} The existing
      \texttt{isls-gateway} (Axum) stays in the workspace. The
      generated Studio is a separate, standalone application.
\item \textbf{Not multi-turn generation.} Still single-shot. Iterative
      refinement is D7+.
\end{itemize}


% ══════════════════════════════════════════════════════════════════════════════
\section{The M\"obius Topology}
% ══════════════════════════════════════════════════════════════════════════════

A M\"obius strip has one continuous surface --- inside and outside are
the same. After D6:

\begin{enumerate}[nosep]
\item \textbf{ISLS} (the generator) produces \textbf{ISLS Studio}
      (the management interface).
\item \textbf{ISLS Studio} displays and manages the \textbf{norms}
      that drive ISLS's generation.
\item Editing norms in the Studio changes what ISLS generates next.
\item What ISLS generates next includes the next version of the Studio.
\end{enumerate}

The generator and the generated share the same data surface. They are
two views of one object --- the M\"obius strip. This is not metaphor.
It is the actual data flow.


% ══════════════════════════════════════════════════════════════════════════════
\section{Architecture Overview}
% ══════════════════════════════════════════════════════════════════════════════

\subsection{Generation Flow}

\begin{lstlisting}[style=bash]
examples/isls_studio.toml
    |
    v
isls forge-v2 --requirements examples/isls_studio.toml
              --api-key $KEY
              --output ./output/isls-studio
    |
    v
HDAG Pipeline (S0-S7):
  48 structural files (models, services, routes, frontend, docker)
  + 7 LLM files (SQL queries for Norm, NormCandidate, etc.)
    |
    v
output/isls-studio/
  backend/    (Rust/Actix-Web, compiles with cargo build)
  frontend/   (Vanilla JS SPA, served by Nginx)
  docker-compose.yml
  .env.example
\end{lstlisting}

\subsection{Entity Model}

The generated application manages these entities:

\begin{longtable}{lp{9cm}}
\toprule
\textbf{Entity} & \textbf{Purpose} \\
\midrule
\texttt{User}
  & JWT authentication (auto-injected by ISLS, always present). \\
\texttt{Norm}
  & A reusable software pattern. Maps to \texttt{isls\_norms::Norm}.
    Fields: name, level, version, trigger keywords, layer count,
    domain count, observation count, builtin flag. \\
\texttt{NormCandidate}
  & A pattern under observation, not yet promoted. Maps to
    \texttt{isls\_norms::NormCandidate}. Fields: status, consistency,
    observation count, domain names, layer count. \\
\texttt{Generation}
  & A record of a forge run. Fields: description, domain, status
    (pending/running/completed/failed), entity count, file count,
    token count, compile ok, duration, output path. \\
\texttt{ScrapeJob}
  & A record of a repository scrape. Fields: source type
    (local/git), source path or URL, domain, status, file count,
    artifact count, language summary. \\
\bottomrule
\end{longtable}

This is a 5-entity application (6 with User). Comparable in complexity
to the D2 E-Commerce test (9 entities) and well within the proven
generation envelope.


% ══════════════════════════════════════════════════════════════════════════════
\section{Work Packages}
% ══════════════════════════════════════════════════════════════════════════════

% ──────────────────────────────────────────────────────────────────────────────
\subsection{W1: Self-Description TOML}
\label{sec:w1}
% ──────────────────────────────────────────────────────────────────────────────

\subsubsection{Goal}

Create \texttt{examples/isls\_studio.toml} --- a standard ISLS
requirements file whose entities describe ISLS's own data types.

\subsubsection{TOML Content}

\begin{lstlisting}[style=toml]
[app]
name = "isls-studio"
description = "ISLS Studio — Self-generated management interface for the norm-driven code generation system"

[backend]
framework = "actix-web"
database = "postgresql"
auth = "jwt"

[frontend]
framework = "vanilla-js"
theme = "dark"

[deployment]
docker = true
ports = { backend = 8080, frontend = 3000, database = 5432 }

# ── Entity: Norm ──────────────────────────────────────────────

[[entities]]
name = "Norm"
fields = [
    { name = "norm_id",           field_type = "String"  },
    { name = "name",              field_type = "String"  },
    { name = "level",             field_type = "String"  },
    { name = "version",           field_type = "String"  },
    { name = "trigger_keywords",  field_type = "String"  },
    { name = "layer_count",       field_type = "i32"     },
    { name = "domain_count",      field_type = "i32"     },
    { name = "observation_count", field_type = "i32"     },
    { name = "is_builtin",        field_type = "bool"    },
    { name = "description",       field_type = "String", nullable = true },
]

# ── Entity: NormCandidate ─────────────────────────────────────

[[entities]]
name = "NormCandidate"
fields = [
    { name = "candidate_id",      field_type = "String"  },
    { name = "status",            field_type = "String"  },
    { name = "consistency",       field_type = "f64"     },
    { name = "observation_count", field_type = "i32"     },
    { name = "domain_names",      field_type = "String"  },
    { name = "layer_count",       field_type = "i32"     },
    { name = "signature",         field_type = "String"  },
]

# ── Entity: Generation ────────────────────────────────────────

[[entities]]
name = "Generation"
fields = [
    { name = "description",   field_type = "String"  },
    { name = "domain",        field_type = "String"  },
    { name = "status",        field_type = "String"  },
    { name = "entity_count",  field_type = "i32"     },
    { name = "file_count",    field_type = "i32"     },
    { name = "token_count",   field_type = "i32"     },
    { name = "compile_ok",    field_type = "bool"    },
    { name = "duration_secs", field_type = "f64"     },
    { name = "output_path",   field_type = "String", nullable = true },
]

# ── Entity: ScrapeJob ─────────────────────────────────────────

[[entities]]
name = "ScrapeJob"
fields = [
    { name = "source_type",      field_type = "String"  },
    { name = "source_path",      field_type = "String"  },
    { name = "domain",           field_type = "String"  },
    { name = "status",           field_type = "String"  },
    { name = "file_count",       field_type = "i32"     },
    { name = "artifact_count",   field_type = "i32"     },
    { name = "language_summary", field_type = "String", nullable = true },
]

# ── Entity: NormWiring ────────────────────────────────────────

[[entities]]
name = "NormWiring"
fields = [
    { name = "norm_a_id",    field_type = "String"  },
    { name = "norm_b_id",    field_type = "String"  },
    { name = "description",  field_type = "String"  },
    { name = "add_services", field_type = "String", nullable = true },
    { name = "add_rules",    field_type = "String", nullable = true },
]
foreign_keys = [
    { target = "Norm" },
]
\end{lstlisting}

\subsubsection{Design Decisions}

\begin{itemize}[nosep]
\item \textbf{Flat fields for complex types.} The real \texttt{Norm} has
      nested \texttt{NormLayers}, \texttt{Vec<TriggerPattern>}, etc.
      These are flattened to strings (\texttt{trigger\_keywords} as
      comma-separated, \texttt{layer\_count} as integer summary).
      This is intentional: the generated Studio is a \emph{dashboard},
      not a runtime. Complex operations happen in the CLI.
\item \textbf{No FK from Generation/ScrapeJob to Norm.} Generations and
      scrapes produce norms indirectly (via the observation pipeline).
      The database link is through domain names, not foreign keys.
\item \textbf{NormWiring has FK to Norm.} This is the only cross-entity
      relationship, and it is genuine (a wiring connects two norms).
\item \textbf{String-typed enums.} \texttt{status}, \texttt{level},
      \texttt{source\_type} are strings, not Rust enums. The LLM
      generates simpler, more reliable SQL for string columns. Enum
      validation is a norm for D7.
\end{itemize}

\subsubsection{W1 Verification}

\begin{enumerate}[nosep]
\item File \texttt{examples/isls\_studio.toml} exists and is valid TOML.
\item \texttt{isls forge-v2 --requirements examples/isls\_studio.toml
      --mock-oracle --output ./output/test-studio} runs without error.
\item Mock output directory contains expected file structure
      (backend/, frontend/, docker-compose.yml).
\end{enumerate}


% ──────────────────────────────────────────────────────────────────────────────
\subsection{W2: Meta-Forge Command}
\label{sec:w2}
% ──────────────────────────────────────────────────────────────────────────────

\subsubsection{Goal}

Add \texttt{isls forge-self} as a dedicated CLI command that generates
ISLS Studio. This is a thin wrapper around \texttt{forge-v2} that:

\begin{enumerate}[nosep]
\item Hardcodes the requirements path to
      \texttt{examples/isls\_studio.toml}.
\item Sets a default output directory of \texttt{./output/isls-studio}.
\item Logs the M\"obius context: ``Generating ISLS Studio --- the
      generator generating itself.''
\item After generation, runs D5's \texttt{ArtifactCollector} on the
      output to self-observe (feeding the D4 norm loop with
      the Studio's own topology).
\end{enumerate}

\subsubsection{Command}

\begin{lstlisting}[style=bash]
isls forge-self --api-key $KEY
isls forge-self --api-key $KEY --output ./custom-path
isls forge-self --mock-oracle    # for testing without API
\end{lstlisting}

\subsubsection{Modified Files}

\begin{longtable}{lp{10cm}}
\toprule
\textbf{File} & \textbf{Change} \\
\midrule
\texttt{isls-cli/src/main.rs}
  & Add \texttt{Command::ForgeSelf} variant. Internally delegates to
    the existing \texttt{forge-v2} codepath with hardcoded TOML path. \\
\bottomrule
\end{longtable}

No new crate, no new module. This is 30--50 lines in \texttt{main.rs}
that reuse the existing \texttt{cmd\_forge\_v2()} logic.

\subsubsection{W2 Verification}

\begin{enumerate}[nosep]
\item \texttt{isls forge-self --mock-oracle} produces output in
      \texttt{./output/isls-studio/}.
\item \texttt{isls forge-self --api-key \$KEY} produces a full
      LLM-generated Studio.
\item The generated \texttt{backend/} compiles with \texttt{cargo build}.
\end{enumerate}


% ──────────────────────────────────────────────────────────────────────────────
\subsection{W3: Builtin Norm Seed}
\label{sec:w3}
% ──────────────────────────────────────────────────────────────────────────────

\subsubsection{Goal}

When ISLS Studio is generated, the database migration should include
seed INSERT statements for the 24 builtin norms. This ensures the
Studio launches with a populated norm catalog, not an empty database.

\subsubsection{Approach}

Extend the structural migration generator (\texttt{structural.rs},
function \texttt{generate\_migration()}) to detect when an entity
named ``Norm'' with field ``is\_builtin'' exists in the AppSpec. If
detected, append INSERT statements for the 24 builtin norms after
the CREATE TABLE statements.

The seed data is derived from \texttt{isls\_norms::builtin\_norms()}
at compile time. Each norm is serialized to its TOML fields:
\texttt{norm\_id}, \texttt{name}, \texttt{level} (as string),
\texttt{version}, \texttt{trigger\_keywords} (joined with ``,``),
\texttt{is\_builtin = true}.

\subsubsection{Constraint}

This is a conditional enhancement to the structural generator --- it
activates ONLY when the entity model contains a ``Norm'' entity with
an \texttt{is\_builtin} field. For all other AppSpecs (warehouse,
hotel, etc.), the migration generator behaves exactly as before.
Rule~8 (seed INSERT only columns from EntityDef) is preserved because
the INSERT columns match the Norm entity's field definitions.

\subsubsection{Modified Files}

\begin{longtable}{lp{10cm}}
\toprule
\textbf{File} & \textbf{Change} \\
\midrule
\texttt{isls-forge-llm/src/structural.rs}
  & In \texttt{generate\_migration()}: after CREATE TABLE for
    ``norms'', if entity has \texttt{is\_builtin} field, append
    INSERT statements from \texttt{builtin\_norms()}. \\
\bottomrule
\end{longtable}

\subsubsection{W3 Verification}

\begin{enumerate}[nosep]
\item Generate Studio with \texttt{forge-self --mock-oracle}.
\item Inspect \texttt{backend/migrations/001\_initial.sql}.
\item Verify 24 INSERT INTO norms statements exist.
\item Verify each INSERT has correct columns matching the Norm entity.
\item Generate a non-ISLS app (e.g. warehouse) --- verify no norm
      seed appears in its migration.
\end{enumerate}


% ──────────────────────────────────────────────────────────────────────────────
\subsection{W4: End-to-End Verification}
\label{sec:w4}
% ──────────────────────────────────────────────────────────────────────────────

\subsubsection{Goal}

Prove that the generated ISLS Studio is a fully operational
application. This is the D6 acceptance test.

\subsubsection{Verification Sequence}

\begin{enumerate}
\item \textbf{Generate:}
      \texttt{isls forge-self --api-key \$KEY}

\item \textbf{Compile:}
      \texttt{cd output/isls-studio/backend \&\& cargo build}

\item \textbf{Deploy:}
      \texttt{cd output/isls-studio \&\& docker-compose up -d}

\item \textbf{Health Check:}
      \texttt{curl http://localhost:8080/api/health}
      --- expect 200 OK.

\item \textbf{Register:}
      \begin{lstlisting}[style=bash]
curl -X POST http://localhost:8080/api/auth/register \
  -H "Content-Type: application/json" \
  -d '{"email":"admin@isls.local","password":"isls2026","role":"admin"}'
      \end{lstlisting}

\item \textbf{Login:}
      \begin{lstlisting}[style=bash]
TOKEN=$(curl -s -X POST http://localhost:8080/api/auth/login \
  -H "Content-Type: application/json" \
  -d '{"email":"admin@isls.local","password":"isls2026"}' \
  | jq -r '.token')
      \end{lstlisting}

\item \textbf{List Norms (seeded):}
      \begin{lstlisting}[style=bash]
curl -H "Authorization: Bearer $TOKEN" \
  http://localhost:8080/api/norms
      \end{lstlisting}
      --- expect 24 builtin norms in response.

\item \textbf{Create Generation Record:}
      \begin{lstlisting}[style=bash]
curl -X POST http://localhost:8080/api/generations \
  -H "Authorization: Bearer $TOKEN" \
  -H "Content-Type: application/json" \
  -d '{"description":"Pet shop","domain":"petshop",
       "status":"completed","entity_count":5,
       "file_count":55,"token_count":12000,
       "compile_ok":true,"duration_secs":45.2}'
      \end{lstlisting}
      --- expect 201 Created.

\item \textbf{Create Scrape Job Record:}
      \begin{lstlisting}[style=bash]
curl -X POST http://localhost:8080/api/scrape_jobs \
  -H "Authorization: Bearer $TOKEN" \
  -H "Content-Type: application/json" \
  -d '{"source_type":"git","source_path":"https://github.com/example/repo",
       "domain":"example","status":"completed",
       "file_count":23,"artifact_count":47}'
      \end{lstlisting}
      --- expect 201 Created.

\item \textbf{Frontend:}
      Open \texttt{http://localhost:3000} in browser.
      Login with admin credentials.
      Navigate to Norms page --- see 24 builtin norms.
      Navigate to Generations page --- see the petshop record.

\item \textbf{Self-Observation:}
      After the Studio is generated and compiled, the D4 observation
      hook has already fed the Studio's topology into the norm
      system. Run \texttt{isls norms candidates} and verify an
      observation from domain ``isls-studio'' exists.
\end{enumerate}


% ══════════════════════════════════════════════════════════════════════════════
\section{File Inventory}
% ══════════════════════════════════════════════════════════════════════════════

\subsection{New Files}

\begin{longtable}{llp{7cm}}
\toprule
\textbf{Location} & \textbf{File} & \textbf{Purpose} \\
\midrule
isls/examples & \texttt{isls\_studio.toml}
  & Self-description TOML \\
\bottomrule
\end{longtable}

\subsection{Modified Files}

\begin{longtable}{llp{7cm}}
\toprule
\textbf{Crate} & \textbf{File} & \textbf{Change} \\
\midrule
isls-cli & \texttt{src/main.rs}
  & Add \texttt{ForgeSelf} command variant (30--50 lines). \\
isls-forge-llm & \texttt{src/structural.rs}
  & Conditional norm seed in migration generator. \\
\bottomrule
\end{longtable}

D6 adds \textbf{1 new file} and modifies \textbf{2 existing files}.
This is intentionally minimal --- the M\"obius inversion is not about
adding code to ISLS. It is about feeding ISLS's own data model through
its existing pipeline.


% ══════════════════════════════════════════════════════════════════════════════
\section{The 12 Rules --- D6 Compliance}
% ══════════════════════════════════════════════════════════════════════════════

\begin{longtable}{clp{9cm}}
\toprule
\textbf{\#} & \textbf{Status} & \textbf{D6 Impact} \\
\midrule
1  & \color{isls-green}\checkmark & The Studio TOML adds no new
     structural/LLM files to the generator. The generated Studio
     follows the standard 48+7 split. \\
2  & \color{isls-green}\checkmark & No module changes. \\
3  & \color{isls-green}\checkmark & Generated CRUD names follow
     convention: \texttt{get\_norm}, \texttt{list\_norms}, etc. \\
4  & \color{isls-green}\checkmark & AuthUser as parameter. \\
5  & \color{isls-green}\checkmark & Owned payloads. \\
6  & \color{isls-green}\checkmark & No actix\_web in queries. \\
7  & \color{isls-green}\checkmark & bcrypt, tracing. \\
8  & \color{isls-green}\checkmark & Norm seed INSERT uses only
     columns from the Norm EntityDef. \\
9  & \color{isls-green}\checkmark & \texttt{rust:1.88-slim}. \\
10 & \color{isls-green}\checkmark & If seed generation fails:
     warn, generate migration without seed. \\
11 & \color{isls-green}\checkmark & ProvidedSymbol unchanged. \\
12 & \color{isls-green}\checkmark & LLM prompts include full
     Norm/NormCandidate/Generation type context. \\
\bottomrule
\end{longtable}


% ══════════════════════════════════════════════════════════════════════════════
\section{Why This Is the M\"obius Moment}
% ══════════════════════════════════════════════════════════════════════════════

The generated Studio is \emph{not} a copy of ISLS. It is ISLS viewed
from the outside --- as a domain model. Consider:

\begin{itemize}[nosep]
\item The ``Norm'' table in the Studio's PostgreSQL is the same data
      that lives in \texttt{\textasciitilde/.isls/norms.json} on the CLI.
\item The ``Generation'' table records what \texttt{isls forge-chat}
      does.
\item The ``ScrapeJob'' table records what \texttt{isls scrape} does.
\item The ``NormCandidate'' table tracks what
      \texttt{observe\_and\_learn()} discovers.
\end{itemize}

The CLI operates on these data structures procedurally. The Studio
operates on them declaratively (CRUD via REST). Same data, two
interfaces, generated by the same pipeline. Inside = Outside.

After D6, the ISLS ecosystem has two views:

\begin{enumerate}[nosep]
\item \textbf{CLI (Gen 1):} Hand-built Rust workspace. 10 crates.
      Procedural pipeline. File-based persistence.
\item \textbf{Studio (Gen 2):} Machine-generated Actix-Web app.
      Declarative CRUD. PostgreSQL persistence.
      Docker-deployable.
\end{enumerate}

D7 will measure whether Gen 2 produces better generation results
than Gen 1 --- by hooking the Studio to the forge pipeline so
that generation requests flow through the web API instead of the
CLI.


% ══════════════════════════════════════════════════════════════════════════════
\section{Token Budget}
% ══════════════════════════════════════════════════════════════════════════════

The Studio has 6 entities (5 + User). Each entity requires $\sim$1 LLM
file (SQL queries). At \textasciitilde\$0.15 per forge run:

\begin{center}
\textbf{D6 generation cost: \textasciitilde\$0.15} (single forge-v2 run)
\end{center}

Additional runs for iteration/fixing: budget \$0.50 total.
From \textasciitilde\$19 remaining: negligible impact.


% ══════════════════════════════════════════════════════════════════════════════
\section{Error Handling}
% ══════════════════════════════════════════════════════════════════════════════

\begin{longtable}{lp{10cm}}
\toprule
\textbf{Failure Mode} & \textbf{Behavior} \\
\midrule
TOML parse error & Standard forge-v2 error handling. \\
LLM query generation fails & S6 Coagula (max 3 cycles). \\
Seed generation fails & Warn, generate migration without seed.
  Studio works but norms table starts empty. \\
\texttt{cargo build} fails & S6 Coagula handles. If 3 cycles
  exhausted: no emission, error logged. \\
Docker-compose fails & Not ISLS's concern --- external
  infrastructure. Report in verification log. \\
\bottomrule
\end{longtable}


% ══════════════════════════════════════════════════════════════════════════════
\section{Constraints for Claude Code}
% ══════════════════════════════════════════════════════════════════════════════

\begin{enumerate}
\item The TOML in \texttt{examples/isls\_studio.toml} MUST use the
      exact same format as \texttt{examples/warehouse.toml}. No new
      TOML keys, no new syntax. The existing parser must handle it
      without modification.
\item \texttt{forge-self} MUST delegate to the existing \texttt{forge-v2}
      codepath. No separate generation pipeline.
\item The norm seed in \texttt{structural.rs} MUST be conditional ---
      it activates only for entities named ``Norm'' with
      \texttt{is\_builtin}. All other AppSpecs are unaffected.
\item The seed INSERTs MUST use columns from the EntityDef (Rule~8).
      The 24 builtin norms are sourced from
      \texttt{isls\_norms::builtin\_norms()} at compile time.
\item Do NOT modify any other structural generators, the HDAG builder,
      the oracle, or the staged closure.
\item Do NOT modify the TOML parser
      (\texttt{isls-hypercube/src/parser.rs}).
\item All new code must compile with \texttt{cargo build --workspace}.
\item Test with \texttt{cargo test --workspace} --- all existing
      tests must pass.
\item Branch: \texttt{claude/implement-isls-d6-spec}.
\item Commit message format: \texttt{D6/W\{n\}: <description>}.
\end{enumerate}


% ══════════════════════════════════════════════════════════════════════════════
\section{Dependency Graph}
% ══════════════════════════════════════════════════════════════════════════════

\begin{verbatim}
W1 (Self-Description TOML)
  |
  +---> W2 (Meta-Forge Command)
  |
  +---> W3 (Builtin Norm Seed)
          |
          v
        W4 (End-to-End Verification)
\end{verbatim}

W1 is the foundation. W2 and W3 can be implemented in parallel
(W2 needs the TOML, W3 modifies the structural generator). W4
requires both W2 and W3 to be complete.

Recommended order: W1 $\to$ W3 $\to$ W2 $\to$ W4.


% ══════════════════════════════════════════════════════════════════════════════
\section{Acceptance Criteria}
\label{sec:acceptance}
% ══════════════════════════════════════════════════════════════════════════════

\begin{longtable}{clp{9cm}}
\toprule
\textbf{\#} & \textbf{Pass} & \textbf{Criterion} \\
\midrule
1 & & \texttt{cargo build --workspace} --- ISLS workspace compiles. \\
2 & & \texttt{cargo test --workspace} --- all tests pass (181+). \\
3 & & \texttt{isls forge-self --mock-oracle} --- produces output
      directory with expected structure. \\
4 & & \texttt{isls forge-self --api-key \$KEY} --- LLM generation
      completes without S6 exhaustion. \\
5 & & Generated \texttt{backend/} compiles with \texttt{cargo build}. \\
6 & & Generated \texttt{migrations/001\_initial.sql} contains
      CREATE TABLE for all 6 entities + 24 norm seed INSERTs. \\
7 & & \texttt{docker-compose up -d} starts all 3 services
      (db, backend, frontend). \\
8 & & Register + login works via REST API. \\
9 & & GET \texttt{/api/norms} returns 24 seeded norms. \\
10 & & CRUD operations on Generation and ScrapeJob work. \\
11 & & Frontend loads in browser, login works, entity pages
       render. \\
12 & & \texttt{isls norms candidates} shows an observation from
       domain ``isls-studio'' (self-observation via D4 hook). \\
\bottomrule
\end{longtable}

\textbf{D6 is passed when criteria 1--6 and 12 are fulfilled.}
Criteria 7--11 (Docker deployment, frontend) are stretch goals
that depend on external infrastructure (Docker, browser) and are
verified manually.


% ══════════════════════════════════════════════════════════════════════════════
\section{Post-D6 Outlook}
% ══════════════════════════════════════════════════════════════════════════════

With D6 complete, the Donnerkeil roadmap enters the \emph{recursive
improvement} phase:

\begin{description}[nosep]
\item[D7 --- Generationsspirale:] Measure Gen~2 (Studio) vs Gen~1
  (CLI). Can the Studio dispatch generation requests that produce
  higher-quality apps? Metrics: compile success rate, token efficiency,
  code quality (measured by D5 topology analysis).
\item[D8 --- Offener Horizont:] Extend beyond CRUD/Actix-Web. Generate
  CLI tools, libraries, data pipelines. PSP-ASIC integration.
  The ``slot machine'' that generates arbitrary infrastructure.
\end{description}

The M\"obius surface is now continuous. Every generation run teaches
the system (D4 observation). External repos feed it (D5 scraping).
The system generates its own management interface (D6). The next
version of the interface will be generated by a system that has
learned from the previous version. The spiral tightens.


% ══════════════════════════════════════════════════════════════════════════════
\section{Verification Sequence (Complete)}
\label{sec:verification}
% ══════════════════════════════════════════════════════════════════════════════

\begin{enumerate}
\item \texttt{cargo build --workspace} --- compiles.
\item \texttt{cargo test --workspace} --- all tests pass.
\item \texttt{isls forge-self --mock-oracle} --- mock output created.
\item Inspect \texttt{output/isls-studio/backend/migrations/001\_initial.sql}
      --- 6 CREATE TABLE + 24 INSERT.
\item \texttt{isls forge-self --api-key \$KEY} --- LLM generation
      completes.
\item \texttt{cd output/isls-studio/backend \&\& cargo build}
      --- compiles.
\item \texttt{isls norms candidates} --- ``isls-studio'' domain
      appears in observations.
\item \texttt{isls norms stats} --- observation count incremented.
\end{enumerate}

\vfill
\begin{center}
\color{isls-gray}\rule{0.5\textwidth}{0.4pt}\\[0.5em]
{\small End of D6 Specification.\\[0.3em]
Der Generator hat sich selbst beschrieben.\\
Jetzt generiert er sich.}
\end{center}

\end{document}
